\section{Introduction}
\label{sec:intro}

Realistic 3D human motions are the “soul’’ of virtual characters, bringing them to life by enabling natural jumping, fighting, and dancing. As portable devices and the metaverse accelerate demand for large-scale character creation, generating adaptive motions from casual video remains challenging. Recent video-to-motion methods like ViMo synthesize diverse 3D motions directly from unconstrained video despite camera motion and occlusion~\cite{vimo2024}, leveraging abundant online footage to produce contextually consistent results. Data-driven generative approaches—such as text-to-motion diffusion models~\cite{motiondiffuse2022} and music-conditioned dance generation like EDGE~\cite{edge2023}—offer scalable alternatives but inherit biases from limited datasets like AMASS~\cite{amass2019}, Human3.6M~\cite{h3wb2014}, and AIST++~\cite{aist2021}. As a result, models trained on narrow domains struggle to generalize to styles such as classical dance or figure skating, reflecting the ongoing scarcity and high cost of collecting diverse 3D motion data.

Casual video offers an escape from this scarcity. Platforms host billions of clips spanning martial arts, ballet, occupational tasks, and impromptu performance. Unlike curated motion-capture datasets, these recordings feature dynamic camera work, shot transitions, occlusion, and varied environments—precisely the failure modes of conventional 3D pose estimators. Methods such as MotionBERT~\cite{motionbert2023} and GLAMR~\cite{glamr2022} regress 3D joints from 2D observations but rely on static or accurately estimable camera parameters. Under casual video conditions, their frame-wise predictions accumulate errors, producing jittery and implausible trajectories. The regression objective itself becomes brittle when perspective shifts violate the single-view assumption.

In our work, we build upon the ViMo framework by integrating diffusion training techniques that enhance both the quality and efficiency of motion generation. Specifically, we adopt the Min-SNR weighting strategy, which treats the denoising process at each timestep as an individual task and assigns loss weights according to the clamped signal-to-noise ratio (SNR)~\cite{minsnr2024}. This strategy balances conflicting gradients across timesteps, resulting in faster convergence and improved motion realism. Furthermore, we incorporate a vectorized timestep formulation inspired by the Probabilistic Timestep Sampling Strategy (PTSS), which enables more flexible and efficient temporal modeling in video diffusion~\cite{ptss2024}. By sampling timesteps probabilistically, our method avoids the combinatorial explosion associated with independent frame timesteps, making training scalable for complex motion sequences.

ViMo~\cite{vimo2024} reframes video-to-motion as a generative task rather than a reconstruction problem. Instead of estimating exact 3D coordinates and camera pose, a diffusion model synthesizes plausible motion sequences conditioned on the 2D pose trajectory extracted from video. Multi-view projections of the generated motion are compared against the input 2D sequence, bypassing explicit camera modeling. The denoising network operates over the full temporal sequence, enforcing coherence through transformer self-attention. This approach tolerates missing frames, occlusions, and rapid viewpoint changes while producing diverse motions that match the stylistic silhouette and rhythm of the source video.

The present work centers on a novel implementation of ViMo enhanced by contemporary diffusion training advances. Two technical improvements define our training regime. First, we integrate the Min-SNR weighting strategy~\cite{minsnr2024}. Standard diffusion training uses uniform loss weighting across noise levels, yet low-SNR steps contribute minimal signal relative to noise, slowing convergence. Min-SNR reweights each timestep by its signal-to-noise ratio, emphasizing high-SNR steps and suppressing low-SNR noise. This yields faster training, lower final loss, and sharper motion details. Second, we adopt the Vectorized Timestep Approach~\cite{ptss2024}, which encodes diffusion timesteps as dense vectors rather than scalar indices. In motion generation, temporal structure matters: the phase relationship between music beats and dance steps, or the cyclic nature of gait, depends on precise timing. Scalar timesteps discard this context. Vectorization embeds timestep information as a function of sequence position, enabling the transformer to exploit long-range temporal correlations. The result is improved rhythmic alignment and smoother interpolation across frames.

Combined, these techniques yield a denoising model that trains efficiently and captures nuanced temporal dynamics while retaining ViMo’s transformer backbone with FiLM conditioning and classifier-free guidance. Our contribution lies not in architectural novelty but in a refined training recipe that delivers robust performance on highly casual video footage. Although this report focuses on the video-to-motion generative core, it forms the first stage of a broader three-stage animation pipeline: extracting kinematic motion from video, translating it into physics-based control policies through imitation learning, and ultimately applying an adaptive refinement loop to correct deployment-time errors. The latter stages remain future work, included here only to situate the present effort within the larger vision.